\section{Publications}
\label{sec:publications}

In this section, we present the student's publications. Below is a list of eight accepted papers where the student had some authorship role\footnote{All accepted publications are directly related to or stem from the research work presented in this proposal, even those where the student is on the record as a co-author.}:

{\footnotesize
\begin{itemize}
    \item \textbf{D. Tadokoro}, R. Siqueira, and P. Meirelles, \textit{Kworkflow: a Linux kernel Developer Automation Workflow System}. In \textit{39th Brazilian Symposium on Software Engineering - Tools Track}.
    \item \textbf{D. Tadokoro} and P. Meirelles, \textit{Can the Linux kernel sustain 30 more years of growth? Toward mitigating bottlenecks in its development model}. In \textit{39th Brazilian Symposium on Software Engineering - Insightful Ideas and Emerging Results Track}.
    \item \textbf{D. Tadokoro}, R. Passos, and P. Meirelles, \textit{Guidelines for Boosting Long-Lasting FLOSS Contributors}. In \textit{DebConf 2025 - Academic Track}.
    \item \textbf{D. Tadokoro}, E. Mendes, and P. Meirelles, \textit{Estudo, implementação e análise de estratégias para simplificar o modelo de contribuição do kernel Linux}. In \textit{32º Simpósio Internacional de Iniciação Científica e Tecnológica da USP Universidade de São Paulo}.
    \item R. Passos, A. Pilone, \textbf{D. Tadokoro}, and P. Meirelles, \textit{Streamlining Analysis on the Linux project with DUKS}. In \textit{VISSOFT 2025 - Visualization Challenge}.
    \item L. Arcanjo, \textbf{D. Tadokoro}, and P. Meirelles, \textit{ArKanjo: a tool for detecting function-level Code Duplication in the Linux Kernel}. In \textit{DebConf 2025 - Academic Track}.
    \item \textbf{D. Tadokoro} and P. Meirelles, \textit{Estratégias de ensino para incentivar a participação consistente em projetos de Software Livre}. In \textit{13th Workshop on Software Visualization, Evolution and Maintenance}.
    \item R. Passos, A. Pilone, \textbf{D. Tadokoro}, and P. Meirelles, \textit{DUKS: visualizações e análises unificadas para o Kernel Linux}. In \textit{13th Workshop on Software Visualization, Evolution and Maintenance}.
\end{itemize}
}
Beyond accepted publications, there are two works in progress to be submitted to the \textit{International Conference on Software Engineering 2026 - Software Engineering Education and Training (ICSE 2026 - SEET)} and the journal \textit{Information and Software Technology (IST)}.
